\setcounter{chapter}{1}
\section*{Devoir maison Noël :  Les Fondamentaux de la Programmation (Variables, Conditions, Boucles)}
\textbf{Consignes : }Pour chaque question,  justifiez vos réponses en expliquant l'exécution du code.


\begin{Exercise}[title={Variables et Opérations}]

\Question Quelle est la valeur finale de la variable c ? Quel est son type ?
\begin{lstlisting}[style=pythonstyle, language=python]
a = 5
b = 2
c = a / b
a = a + 1
c = c * a
\end{lstlisting}


\Question Après l'exécution du code ci-dessous, quelle est la valeur finale de la variable y ? Quel est son type ?
\begin{lstlisting}[style=pythonstyle, language=python]
x = 5
y = 2
x = (x + 1) * y
y = x % 5  # Opérateur modulo
\end{lstlisting}
\end{Exercise}

\begin{Exercise}[title={Conditions (If / Else)}]
\Question Soient $A=10$ et $B=5$. Évaluez l'expression booléenne suivante en tenant compte de la priorité des opérateurs logiques (ET avant OU). Cette expression renvoie-t-elle VRAI, FAUX, une erreur de syntaxe ou dépend-elle du langage ?

$$(A < B) \text{ OU } ((A + B) > 12) \text{ ET } (B \ge 5)$$

\Question Quelle sera la valeur finale de la variable résultat si la variable age vaut 17 ?
\begin{lstlisting}[style=pythonstyle, language=python]
if age < 13:
    resultat = "Enfant"
elif age >= 18:
    resultat = "Adulte"
else:
    resultat = "Adolescent"
\end{lstlisting}


\Question Quel sera le message affiché par ce code avec les valeurs initiales $\text{temp} = 25$ et $\text{pluie} = \text{Vrai}$ (True)?
\begin{lstlisting}[style=pythonstyle, language=python]
if temp > 20:
    if pluie:
        message = "Journée fraîche et pluvieuse"
    else:
        message = "Journée ensoleillée"
else:
    message = "Journée froide"
\end{lstlisting}

\Question Quelle est la condition logique qui vérifie si la variable x N'EST PAS comprise dans l'intervalle fermé $[10, 20]$ (c'est-à-dire si $x < 10$ ou $x > 20$)?
\begin{enumerate}
    \item x > 10 ET x < 20
    \item x >= 10 OU x <= 20
    \item x < 10 OU x > 20
    \item NON (x >= 10 ET x <= 20)
\end{enumerate}
\end{Exercise}


\begin{Exercise}[title={Boucles (For / While)}]
\Question Quelle est la valeur finale de la variable somme après l'exécution de cette boucle?

\begin{lstlisting}[style=pythonstyle, language=python]
somme = 0
# range(1, 5) génère les nombres 1, 2, 3, 4
for i in range(1, 5):
    somme = somme + i
\end{lstlisting}

\Question Combien de fois le corps de la boucle while est-il exécuté?
\begin{lstlisting}[style=pythonstyle, language=python]
compteur = 0
a = 1
while a <= 5:
    a = a + 2
    compteur = compteur + 1
\end{lstlisting}

\Question Quel bloc de code placé à l'intérieur de la boucle pourrait donner une boucle infinie?
\begin{lstlisting}[style=pythonstyle, language=python]
i = 10
while i > 0:
    # BLOC DE CODE
\end{lstlisting}

\Question Combien de fois le message "OK" sera-t-il affiché par ce bloc de code ?
\begin{lstlisting}[style=pythonstyle, language=python]
# range(1, 11) génère les nombres 1, 2, ..., 10
for k in range(1, 11):
    if k % 3 == 0:
        print("OK")
\end{lstlisting}
\end{Exercise}



\begin{Exercise}[title={En autonomie (ou presque) : Calculateur de Points Séquentiel}]
Une équipe de sport joue une série de 5 matchs. Nous allons calculer son score total en demandant le résultat de chaque match à l'utilisateur, en respectant les règles suivantes :
\begin{itemize}
    \item Victoire (V) : +3 points
    \item Nul (N) : +1 point
    \item Défaite (D) : +0 point
\end{itemize}
Règle du Bonus spécial : si l'équipe a marqué au moins 7 points à la fin de la série, elle reçoit un point bonus.

\subsubsection*{Etape 1 : Initialisation des variables (Variables)}
\begin{enumerate}
    \item Déclarez une variable nommée score\_total et initialisez-la à 0.
    \item Déclarez une variable de contrôle de la boucle nommée match\_actuel et initialisez-la à 1 (pour compter le match que nous sommes en train de traiter).
\end{enumerate}

\subsubsection*{Etape 2 : Boucle conditionnelle (Boucles)}
Utilisez une boucle while pour vous assurer que nous traitons exactement 5 matchs. La boucle doit continuer tant que le match\_actuel est inférieur ou égal à 5.

\subsubsection*{Étape 3 : Saisie, Conditions et Incrémentation (Variables, Conditions)}
À l'intérieur de la boucle, vous devez :
\begin{enumerate}
    \item Demander à l'utilisateur d'entrer le résultat (V, N, ou D) pour le match\_actuel.
    \item Utiliser les conditions (if, elif, else) pour mettre à jour le score\_total.
    \item Important : Incrémenter la variable match\_actuel pour éviter la boucle infinie et passer au match suivant.
\end{enumerate}
\vspace{-1cm}
\subsubsection*{Étape 4 : Appliquer le bonus (Condition Finale)}
Une fois la boucle while terminée, vous appliquez la règle du bonus exactement comme dans l'exercice précédent.
\end{Exercise}


\begin{Exercise}[title={Conversions de Bases Numériques}]
\subsection*{Partie A : Binaire (Base 2) $\leftrightarrow$ Décimal (Base 10)}
\Question Convertir les nombres entre la base 2 et la base 10. JUSTIFIEZ VOS REPONSES.\\[0.2cm]
\begin{tabular}{|l|l|l|l|l|}
    \hline
        Question & Nombre à convertir & Base de départ & Base d'arrivée & Réponse \\ \hline
        A1 & ${10110}_{(2)}$ & Binaire & Décimal & ~ \\ \hline
        A2 & ${111111}_{(2)}$ & Binaire & Décimal & ~ \\ \hline
        A3 & $45_{(10)}$ & Décimal & Binaire & ~ \\ \hline
        A4 & $100_{(10)}$ & Décimal & Binaire & ~ \\ \hline
\end{tabular}

%
\subsection*{Partie B : Binaire (Base 2) $\leftrightarrow$ Hexadécimal (Base 16)}

\Question Convertir les nombres entre la base 2 et la base 16, en utilisant le groupement par paquets de 4 bits. JUSTIFIEZ VOS REPONSES.\\

Rappel : $A=10, B=11, C=12, D=13, E=14, F=15$.\\

\begin{tabular}{|l|l|l|l|l|}
    \hline
        Question & Nombre à convertir & Base de départ & Base d'arrivée & Réponse \\ \hline
        B1 & ${10101101}_{(2)}$ & Binaire & Hexadécimal & ~ \\ \hline
        B2 & ${001101011110}_{(2)}$ & Binaire & Hexadécimal & ~ \\ \hline
        B3 & $9C_{16}$ & Hexadécimal & Binaire & ~ \\ \hline
        B4 & $1F8_{16}$ & Hexadécimal & Binaire & ~ \\ \hline
\end{tabular}


\subsection*{Partie C : Hexadécimal (Base 16) $\leftrightarrow$ Décimal (Base 10)}

\Question Convertir les nombres entre la base 16 et la base 10. JUSTIFIEZ VOS REPONSES. \\[0.2cm]
\begin{tabular}{|l|l|l|l|l|}
    \hline
        Question & Nombre à convertir & Base de départ & Base d'arrivée & Réponse \\ \hline
        C1 & $2B_{16}$ & Hexadécimal & Décimal & ~ \\ \hline
        C2 & $FF_{16}$ & Hexadécimal & Décimal & ~ \\ \hline
        C3 & $78_{10}$ & Décimal & Hexadécimal & ~ \\ \hline
        C4 & $300_{10}$  & Décimal & Hexadécimal & ~ \\ \hline
\end{tabular}

\end{Exercise}